\chapter{Pruebas y Validación}

El presente capítulo describe la estrategia de validación implementada para garantizar la corrección funcional del algoritmo y evaluar el comportamiento de las diferentes estrategias de paralelización desarrolladas durante el proyecto.

Las pruebas realizadas tuvieron como objetivo verificar que todas las implementaciones produjeran resultados consistentes respecto a la línea base secuencial y que las mejoras en desempeño no afectaran la calidad de la solución encontrada.

Asimismo, se analizaron métricas clásicas de computación paralela tales como speedup, eficiencia y escalabilidad, complementadas con la Ley de Amdahl para estimar el potencial máximo de aceleración del problema.

\section{Metodología de Validación}

La validación se desarrolló siguiendo las siguientes etapas:

\begin{enumerate}
    \item Verificación funcional de la implementación secuencial.
    \item Comparación de resultados entre implementaciones.
    \item Validación de restricciones matemáticas.
    \item Verificación de consistencia mediante ROC-AUC.
    \item Evaluación del desempeño computacional.
    \item Medición de speedup.
    \item Cálculo de eficiencia.
    \item Estudio de escalabilidad.
\end{enumerate}

\section{Entorno Experimental}

Durante la ejecución de las pruebas debe registrarse la configuración del entorno utilizado.

\begin{table}[H]
\centering
\caption{Configuración del entorno experimental}
\begin{tabular}{ll}
\toprule
Característica & Especificación \\
\midrule
Sistema Operativo  & Windows 11 (64 bits) \\
Procesador         & Intel Core i5 \\
Núcleos físicos     & 4 \\
Hilos lógicos       & 8 \\
Memoria RAM         & 8 GB \\
GPU                 & NVIDIA (Google Colab -- Tesla T4) \\
Versión Python      & 3.11 \\
Versión GCC         & 11.4.0 (MinGW-w64 / WSL) \\
Versión OpenMP      & 4.5 \\
Versión MPI         & MPICH 4.1 \\
Versión CUDA        & 12.2 \\
Versión PyCUDA      & 2024.1 \\
\bottomrule
\end{tabular}
\end{table}

Los resultados obtenidos deben almacenarse en el archivo:

\begin{center}
\texttt{Resultados/benchmark.csv}
\end{center}

Este archivo centraliza las métricas de desempeño generadas por las diferentes implementaciones y facilita la construcción de tablas y gráficas comparativas.

%%%%%%%%%%%%%%%%%%%%%%%%%%%%%%%%%%%%%%%%%%%%%%%%%%%%%%
\section{Pruebas Funcionales}
%%%%%%%%%%%%%%%%%%%%%%%%%%%%%%%%%%%%%%%%%%%%%%%%%%%%%%

Las pruebas funcionales verifican que todas las implementaciones produzcan resultados correctos desde el punto de vista matemático.

\subsection{Prueba PF-01: Correctitud Secuencial}

\textbf{Objetivo}

Verificar que la implementación secuencial produzca resultados válidos.

\textbf{Procedimiento}

\begin{enumerate}
    \item Ejecutar la versión secuencial.
    \item Registrar el mejor ROC-AUC encontrado.
    \item Registrar el vector de pesos asociado.
    \item Verificar el cumplimiento de las restricciones del simplex.
\end{enumerate}

\textbf{Resultado esperado}

\begin{itemize}
    \item ROC-AUC válido.
    \item Restricciones satisfechas.
    \item Solución reproducible utilizando la misma semilla aleatoria.
\end{itemize}

\textbf{Resultado obtenido}

Completar durante la experimentación.

%%%%%%%%%%%%%%%%%%%%%%%%%%%%%%%%%%%%%%%%%%%%%%%%%%%%%%

\subsection{Prueba PF-02: Equivalencia Multicore}

\textbf{Objetivo}

Comprobar que la implementación multicore produzca resultados equivalentes a la línea base.

\[
AUC_{multicore}\approx AUC_{secuencial}
\]

La diferencia observada debe atribuirse únicamente al carácter aleatorio de la exploración y no a errores computacionales.

%%%%%%%%%%%%%%%%%%%%%%%%%%%%%%%%%%%%%%%%%%%%%%%%%%%%%%

\subsection{Prueba PF-03: Equivalencia OpenMP}

\[
AUC_{OpenMP}\approx AUC_{secuencial}
\]

El mejor resultado encontrado debe mantenerse dentro del mismo rango de calidad obtenido por la implementación secuencial.

%%%%%%%%%%%%%%%%%%%%%%%%%%%%%%%%%%%%%%%%%%%%%%%%%%%%%%

\subsection{Prueba PF-04: Equivalencia MPI}

\[
AUC_{MPI}\approx AUC_{secuencial}
\]

Se verifica que la distribución de candidatos entre procesos no altere la validez matemática de la solución.

%%%%%%%%%%%%%%%%%%%%%%%%%%%%%%%%%%%%%%%%%%%%%%%%%%%%%%

\subsection{Prueba PF-05: Equivalencia CUDA}

\[
AUC_{CUDA}\approx AUC_{secuencial}
\]

La aceleración proporcionada por la GPU no debe modificar el resultado final del algoritmo.

%%%%%%%%%%%%%%%%%%%%%%%%%%%%%%%%%%%%%%%%%%%%%%%%%%%%%%

\subsection{Prueba PF-06: Restricciones del Simplex}

\textbf{Objetivo}

Validar que todos los vectores de pesos generados cumplan:

\[
\sum_{i=1}^{N}W_i=1
\]

y

\[
W_i \ge 0
\]

\textbf{Resultado esperado}

Todos los candidatos deben satisfacer las restricciones del simplex.

%%%%%%%%%%%%%%%%%%%%%%%%%%%%%%%%%%%%%%%%%%%%%%%%%%%%%%

\subsection{Prueba PF-07: Consistencia del Scoring}

Según lo establecido en el Capítulo 3, una solución es considerada consistente cuando:

\[
ROC\text{-}AUC \ge 0.8
\]

Esta prueba verifica que los mejores candidatos encontrados alcancen niveles adecuados de discriminación entre las clases analizadas.

%%%%%%%%%%%%%%%%%%%%%%%%%%%%%%%%%%%%%%%%%%%%%%%%%%%%%%
\section{Pruebas de Desempeño}
%%%%%%%%%%%%%%%%%%%%%%%%%%%%%%%%%%%%%%%%%%%%%%%%%%%%%%

Las pruebas de desempeño buscan cuantificar el impacto de cada estrategia de paralelización.

Cada experimento debe ejecutarse múltiples veces para reducir la influencia de factores externos como carga del sistema operativo, procesos concurrentes o variaciones en la frecuencia del procesador.

Los experimentos se realizaron utilizando \texttt{k\_candidates = 1{,}000{,}000} y 8 workers para las implementaciones paralelas basadas en CPU (Multicore, OpenMP y MPI), mientras que la implementación CUDA utilizó 1 GPU.

\subsection{Tiempo de ejecución}

\begin{table}[H]
\centering
\caption{Tiempo promedio de ejecución}
\begin{tabular}{lccc}
\toprule
Implementación & Corrida 1 (s) & Corrida 2 (s) & Promedio (s) \\
\midrule
Secuencial & 830.2724 & 801.3928 & \textbf{815.8326} \\
Multicore  & 173.7440 & 177.7167 & \textbf{175.7304} \\
OpenMP     & 0.024036 & 0.056817 & \textbf{0.040427} \\
MPI        & 0.027697 & 0.070864 & \textbf{0.049281} \\
CUDA       & 0.015747 & 0.039362 & \textbf{0.027555} \\
\bottomrule
\end{tabular}
\end{table}

\textit{Nota: Multicore, OpenMP y MPI ejecutados con 8 workers; CUDA ejecutado sobre 1 GPU.}

%%%%%%%%%%%%%%%%%%%%%%%%%%%%%%%%%%%%%%%%%%%%%%%%%%%%%%
\section{Speedup}
%%%%%%%%%%%%%%%%%%%%%%%%%%%%%%%%%%%%%%%%%%%%%%%%%%%%%%

El speedup mide la aceleración obtenida respecto a la versión secuencial.

\begin{equation}
S_p=
\frac{T_{seq}}
     {T_p}
\end{equation}

donde:

\begin{itemize}
    \item $T_{seq}$ es el tiempo secuencial.
    \item $T_p$ es el tiempo paralelo.
\end{itemize}

Utilizando $T_{seq} = 815.8326$ s como línea base:

\begin{table}[H]
\centering
\caption{Speedup obtenido}
\begin{tabular}{lc}
\toprule
Implementación & Speedup \\
\midrule
Multicore & 4.64$\times$ \\
OpenMP    & 20{,}182$\times$ \\
MPI       & 16{,}557$\times$ \\
CUDA      & 29{,}612$\times$ \\
\bottomrule
\end{tabular}
\end{table}

\textbf{Observaciones clave:}

\begin{itemize}
    \item Multicore obtiene un speedup de aproximadamente $4.6\times$ frente a Python secuencial utilizando 8 procesos, resultado limitado por el GIL de Python y el overhead asociado a la creación de procesos.
    \item OpenMP y MPI, al estar implementados en C, alcanzan aceleraciones de entre $16{,}000\times$ y $20{,}000\times$ respecto a la versión secuencial en Python, reflejando tanto la diferencia de lenguaje como el efecto de la paralelización.
    \item CUDA constituye la implementación más rápida, con un speedup cercano a $29{,}600\times$, aprovechando el paralelismo masivo disponible en la GPU.
\end{itemize}

\textit{Nota: dado que la línea base corresponde a Python (lenguaje interpretado) mientras que OpenMP, MPI y CUDA están implementados en C/CUDA (lenguajes compilados), el speedup reportado combina el efecto de la paralelización con el efecto del cambio de lenguaje. Por esta razón, los valores de eficiencia y la estimación de la fracción paralelizable mediante la Ley de Amdahl (Secciones \ref{sec:eficiencia} y \ref{sec:amdahl_empirico}) no deben interpretarse en el sentido estrictamente clásico, y se recomienda complementarlas con una comparación adicional entre implementaciones del mismo lenguaje (p. ej. OpenMP vs. MPI) para un análisis de escalabilidad más riguroso.}

%%%%%%%%%%%%%%%%%%%%%%%%%%%%%%%%%%%%%%%%%%%%%%%%%%%%%%
\section{Eficiencia}
\label{sec:eficiencia}
%%%%%%%%%%%%%%%%%%%%%%%%%%%%%%%%%%%%%%%%%%%%%%%%%%%%%%

La eficiencia evalúa qué tan bien son aprovechados los recursos paralelos disponibles.

\begin{equation}
E_p=
\frac{S_p}{P}
\end{equation}

donde:

\begin{itemize}
    \item $S_p$ es el speedup observado.
    \item $P$ representa el número de unidades de ejecución.
\end{itemize}

\begin{table}[H]
\centering
\caption{Eficiencia obtenida}
\begin{tabular}{lcc}
\toprule
Implementación & Procesadores ($P$) & Eficiencia \\
\midrule
Multicore & 8 & 0.58 \\
OpenMP    & 8 & 2{,}522.75 \\
MPI       & 8 & 2{,}069.63 \\
CUDA      & 1 (GPU) & 29{,}612.00 \\
\bottomrule
\end{tabular}
\end{table}

\textit{Nota: los valores de eficiencia superiores a 1 obtenidos para OpenMP, MPI y CUDA son consecuencia de comparar implementaciones en C/CUDA contra una línea base en Python, y no deben interpretarse como eficiencia paralela clásica (la cual está acotada entre 0 y 1). Solo el valor obtenido para Multicore, al compartir el mismo lenguaje que la línea base, es directamente comparable bajo la definición clásica de eficiencia.}

%%%%%%%%%%%%%%%%%%%%%%%%%%%%%%%%%%%%%%%%%%%%%%%%%%%%%%
\section{Escalabilidad}
%%%%%%%%%%%%%%%%%%%%%%%%%%%%%%%%%%%%%%%%%%%%%%%%%%%%%%

La escalabilidad mide la capacidad del algoritmo para aprovechar recursos adicionales a medida que se incrementa el número de procesadores.

\begin{table}[H]
\centering
\caption{Escalabilidad}
\begin{tabular}{cccc}
\toprule
Procesadores & Tiempo (s) & Speedup & Eficiencia \\
\midrule
1  & \multicolumn{3}{c}{Sin datos disponibles} \\
2  & \multicolumn{3}{c}{Sin datos disponibles} \\
4  & \multicolumn{3}{c}{Sin datos disponibles} \\
8  & 0.040427 (OpenMP) & 20{,}182$\times$ & 2{,}522.75 \\
16 & \multicolumn{3}{c}{Sin datos disponibles} \\
\bottomrule
\end{tabular}
\end{table}

\textit{Nota: únicamente se dispone de mediciones experimentales para $P=8$. El estudio de escalabilidad completo (variando el número de procesadores entre 1 y 16) queda pendiente de ejecución y debe completarse con corridas adicionales antes de la entrega final.}

%%%%%%%%%%%%%%%%%%%%%%%%%%%%%%%%%%%%%%%%%%%%%%%%%%%%%%
\section{Ley de Amdahl}
%%%%%%%%%%%%%%%%%%%%%%%%%%%%%%%%%%%%%%%%%%%%%%%%%%%%%%

La Ley de Amdahl establece un límite teórico para la aceleración máxima alcanzable por una aplicación paralela.

Sea \(f\) la fracción paralelizable del algoritmo y \(P\) el número de procesadores utilizados.

El speedup teórico viene dado por:

\begin{equation}
S(P)=
\frac{1}
     {(1-f)+\frac{f}{P}}
\end{equation}

Cuando el número de procesadores tiende a infinito:

\[
\lim_{P\to\infty}S(P)
=
\frac{1}{1-f}
\]

Esto implica que incluso disponiendo de recursos ilimitados, la parte secuencial del algoritmo impone una cota superior al rendimiento alcanzable.

En este proyecto, la evaluación independiente de candidatos permite esperar valores elevados de \(f\), ya que la mayor parte del tiempo de ejecución corresponde a tareas paralelizables.

%%%%%%%%%%%%%%%%%%%%%%%%%%%%%%%%%%%%%%%%%%%%%%%%%%%%%%
\section{Estimación Empírica de la Fracción Paralelizable}
\label{sec:amdahl_empirico}
%%%%%%%%%%%%%%%%%%%%%%%%%%%%%%%%%%%%%%%%%%%%%%%%%%%%%%

A partir del speedup experimental observado puede estimarse la fracción paralelizable utilizando:

\[
f=
\frac{1-\frac{1}{S(P)}}
     {1-\frac{1}{P}}
\]

Esta estimación permite cuantificar qué porcentaje del algoritmo realmente se beneficia de la paralelización.

\begin{table}[H]
\centering
\caption{Estimación empírica de la fracción paralelizable}
\begin{tabular}{lccc}
\toprule
Implementación & P & Speedup & f \\
\midrule
Multicore & 8 & 4.64$\times$  & 0.90 \\
OpenMP    & 8 & 20{,}182$\times$ & N/D$^{*}$ \\
MPI       & 8 & 16{,}557$\times$ & N/D$^{*}$ \\
CUDA      & 1 & 29{,}612$\times$ & N/D$^{*}$ \\
\bottomrule
\end{tabular}
\end{table}

\textit{$^{*}$N/D: la fórmula arroja valores de $f$ superiores a 1 (matemáticamente inconsistentes con el modelo de Amdahl) debido a que el speedup involucrado mezcla el efecto de la paralelización con el cambio de lenguaje (Python a C/CUDA). Para CUDA, adicionalmente, $P=1$ genera una división por cero en el denominador. Una estimación válida de $f$ requeriría comparar cada implementación contra una línea base secuencial escrita en el mismo lenguaje.}

Únicamente el valor obtenido para Multicore resulta consistente con el modelo teórico, indicando que aproximadamente el 90\% del algoritmo es paralelizable bajo esa comparación.

%%%%%%%%%%%%%%%%%%%%%%%%%%%%%%%%%%%%%%%%%%%%%%%%%%%%%%
\section{Análisis de Resultados}
%%%%%%%%%%%%%%%%%%%%%%%%%%%%%%%%%%%%%%%%%%%%%%%%%%%%%%

A partir de las tablas anteriores se identifican los siguientes hallazgos:

\begin{itemize}
    \item Todas las implementaciones paralelas superan ampliamente a la versión secuencial en tiempo de ejecución, sin comprometer la calidad de la solución encontrada.
    \item El overhead de comunicación afecta de forma más notoria a Multicore, dado que la creación de procesos y la serialización de datos representan un costo significativo frente al tiempo total de cómputo.
    \item OpenMP presenta un desempeño levemente superior a MPI en el entorno de un solo nodo evaluado, consistente con su menor overhead de sincronización frente al paso de mensajes.
    \item CUDA ofrece el mayor beneficio absoluto gracias al paralelismo masivo de la GPU, aunque su ventaja relativa frente a OpenMP y MPI es moderada una vez que ambos se ejecutan en C.
    \item La comparación de speedup y eficiencia entre implementaciones de distinto lenguaje (Python vs. C/CUDA) debe interpretarse con cautela, ya que combina el efecto de la paralelización con el del cambio de lenguaje.
    \item La estimación empírica de la fracción paralelizable solo resultó matemáticamente consistente para Multicore, evidenciando la necesidad de una línea base secuencial en C para evaluar correctamente OpenMP, MPI y CUDA bajo la Ley de Amdahl.
    \item Como limitación principal se identifica la ausencia de un estudio de escalabilidad con múltiples valores de $P$ (1, 2, 4, 8, 16), pendiente de ejecución experimental.
\end{itemize}

Los hallazgos obtenidos servirán como fundamento para la interpretación de resultados experimentales y las conclusiones finales del proyecto.